% Physical constants and notation used in simulation and learning experiments.
% Numeric values match backend/environment_definition/constants/ at report revision;
% see docs/presentation/ for the live parameter record.

\subsection{Abbreviations}
\begin{tabular}{@{}p{0.12\linewidth}p{0.80\linewidth}@{}}
  \toprule
  Abbreviation & Definition \\
  \midrule
  CNN & Convolutional neural network (vision observation encoder) \\
  EO & Earth observation \\
  FOV & Field of view \\
  GSD & Ground sample distance \\
  LLA & Latitude, longitude, altitude (geodetic coordinates) \\
  LOS & Line of sight \\
  MLP & Multilayer perceptron (scalar observation branch) \\
  MPO & Maximum \emph{a posteriori} policy optimization \\
  OBC & On-board computer \\
  PD & Proportional--derivative (pointing controller) \\
  RL & Reinforcement learning \\
  RW & Reaction wheel \\
  WGS84 & World Geodetic System 1984 (reference ellipsoid) \\
  \bottomrule
\end{tabular}

\subsection{Symbols}
\begin{tabular}{@{}p{0.14\linewidth}p{0.78\linewidth}@{}}
  \toprule
  Symbol & Meaning \\
  \midrule
  $R_{\mathrm{E}}$ & Mean Earth radius (orbital geometry) \\
  $\mu$ & Earth gravitational parameter \\
  $h$ & Orbit altitude above mean Earth radius \\
  $I_{\mathrm{sat}}$ & Spacecraft moment of inertia (2D model: $I_{zz}$) \\
  $\boldsymbol{\tau}$ & Reaction-wheel torque command [N\,m] (scalar in the 2D kernel) \\
  $\theta_{\mathrm{orb}}$ & Orbit-plane angle of the subsatellite track \\
  $\theta_{\mathrm{body}}$ & Body $z$-axis angle in the orbit disk \\
  $\omega_{\mathrm{sat}}$, $\omega_{\mathrm{w}}$ & Spacecraft body rate and wheel speed \\
  $f_{\mathrm{cloud}}$ & Primary-camera cloud-blocked fraction along the sensor strip \\
  \bottomrule
\end{tabular}

\subsection{Physical and simulation constants}
Unless noted, values are fixed inputs; $h$ and cloud layouts are randomized per episode (Table~\ref{tab:constants-mission}).

\begin{table}[htbp]
  \centering
  \caption{Earth, orbit, and spacecraft parameters.}
  \label{tab:constants-earth-spacecraft}
  \small
  \begin{tabular}{@{}llS[table-format=4.4]l@{}}
    \toprule
    {Quantity} & {Symbol} & {Value} & {Unit / note} \\
    \midrule
    Mean Earth radius & $R_{\mathrm{E}}$ & 6371.0088 & \si{\kilo\meter} \\
    Gravitational parameter & $\mu$ & 398600.4418 & \si{\kilo\meter\cubed\per\second\squared} \\
    Geodetic model & --- & \multicolumn{2}{l}{WGS84 ellipsoid (surface / LOS)} \\
    Orbit altitude (episode) & $h$ & \multicolumn{2}{l}{uniform on $[\SI{510}{\kilo\meter},\,\SI{570}{\kilo\meter}]$} \\
    Satellite mass & $m_{\mathrm{sat}}$ & 250 & \si{\kilogram} \\
    Moment of inertia (3D) & $I_{xx}, I_{yy}, I_{zz}$ & \multicolumn{2}{l}{16.6, 21.7, 31.2 \si{\kilogram\meter\squared}} \\
    Moment of inertia (2D sim.) & $I_{\mathrm{sat}}$ & 31.2 & \si{\kilogram\meter\squared} ($I_{zz}$) \\
    RW max torque & $\tau_{\max}$ & 0.1 & \si{\newton\meter} \\
    RW max momentum & $H_{\max}$ & 0.4 & \si{\newton\meter\second} \\
    Star-tracker rate limit & --- & 3 & \si{\degree\per\second} \\
    Off-nadir hard limit & --- & 45 & \si{\degree} \\
    \bottomrule
  \end{tabular}
  \\[0.4em]
  {\footnotesize Sources: \texttt{constants/EARTH.py}, \texttt{constants/MISSION.py}, \texttt{constants/SATELLITE.py}, \texttt{constants/ATTITUDE\_SAFETY.py}.}
\end{table}

\begin{table}[htbp]
  \centering
  \caption{Camera, simulation timing, and sensing discretization.}
  \label{tab:constants-camera-sim}
  \small
  \begin{tabular}{@{}llS[table-format=4.0]l@{}}
    \toprule
    {Quantity} & {Symbol} & {Value} & {Unit / note} \\
    \midrule
    Sensor resolution & --- & \multicolumn{2}{l}{9344 $\times$ 7000 pixels (cross $\times$ along)} \\
    Pixel pitch & --- & 3.2 & \si{\micro\meter} \\
    Focal length & $f$ & 1067 & \si{\milli\meter} \\
    Exposure time & $t_{\mathrm{exp}}$ & 100 & \si{\micro\second} \\
    Strip ray samples & --- & 96 & per primary-camera frame \\
    Observation-line bins (primary) & --- & 100 & vertical FOV discretization \\
    Simulation step & $\Delta t$ & 0.4 & \si{\second} \\
    Episode step cap & --- & 1000 & steps (\texttt{max\_episode\_steps}) \\
    Max captures per orbit & --- & 10 & OBC memory / downlink budget \\
  \bottomrule
  \end{tabular}
  \\[0.4em]
  {\footnotesize Sources: \texttt{constants/SATELLITE.py}, \texttt{constants/SIMULATION.py}.}
\end{table}

\begin{table}[htbp]
  \centering
  \caption{Mission profile and reward scales (notebook~07 / 08 defaults).}
  \label{tab:constants-mission}
  \small
  \begin{tabular}{@{}p{0.36\linewidth}p{0.54\linewidth}@{}}
    \toprule
    Quantity & Value / range \\
    \midrule
    Target grid (baseline) & 50 areas, \SI{15}{\kilo\meter} extent, \SI{40}{\kilo\meter} spacing along meridian \\
    Target anchor latitude & \SI{80.0}{\degree}N (grid start) \\
    Cloud patches per episode & uniform integer in $[15,\,30]$ along corridor \\
    Cloud horizontal extent & \SIrange{1}{80}{\kilo\meter} \\
    Cloud base altitude & \SIrange{4}{12}{\kilo\meter} \\
    Cloud thickness & \SIrange{1}{16}{\kilo\meter}; top capped at \SI{20}{\kilo\meter} \\
    Capture reward weight & $w_{\mathrm{capture}} = 100$ \\
    Optimal / threshold ground range & $d_{\mathrm{op}} = \SI{400}{\kilo\meter}$, $d_{\mathrm{th}} = \SI{1500}{\kilo\meter}$ \\
    Viewing-distance outer gate & \SI{2500}{\kilo\meter} \\
    \bottomrule
  \end{tabular}
  \\[0.4em]
  {\footnotesize Sources: \texttt{s01\_utils/baseline\_overflight.py}, \texttt{constants/AUTONOMOUS\_CONTROL\_REWARD.py}, \texttt{constants/MISSION.py}.}
\end{table}
